% Provenance (never printed): see sources/outline_sources.md (Year 2: radiative forcing
% quantified; feedbacks and climate sensitivity; detection and attribution; scenarios; the
% remaining carbon budget; components of sea-level rise; physical science only). Every statement
% from IPCC AR6 WG1 (SPM; chapters 3, 5, 7 and 9; Annex III data), ledger rows in
% sources/ledger/book2.md (ipcc:*) and sources/ledger/book4.md (ar6*:*). Curves and numbers from
% figdata/bachelor-2/anthropogenic.py and the Book 4 second-writer scratch calculator calc/ch27.py.
\chapter{Changement climatique anthropique}\label{ch:b2:anthropogenic-climate-change}

Depuis 1970, la surface du globe s'est réchauffée plus vite que pendant toute autre période de cinquante ans
des deux derniers millénaires au moins (\emph{confiance élevée}). Les volumes scolaires ont dit ce qui a changé
et pourquoi. Ce chapitre transforme ce récit en bilan~: combien de watts par mètre carré les facteurs humains
ajoutent, combien la planète plus chaude en renvoie vers l'espace, combien de réchauffement chaque tonne de
dioxyde de carbone procure, et comment la hausse du niveau de la mer se partage entre l'eau qui se réchauffe et
la glace qui fond. Chaque nombre vient avec son incertitude, dans le langage calibré des évaluations.
% ledger: ipcc:fastest-2000yr, ipcc:likelihood

\begin{recall}
\textbf{Acquis des volumes scolaires}~: forçage radiatif, $\qty{3.93}{W.m^{-2}}$ pour un doublement du dioxyde
de carbone, et sa loi logarithmique~; rétroaction climatique~; \omterm{def:b2:anthropogenic-climate-change:ecs}{sensibilité climatique à l'équilibre}, meilleure
estimation \qty{3}{\celsius}, \emph{probable} de 2.5 à \qty{4}{\celsius}~; \omterm{def:b2:anthropogenic-climate-change:da}{attribution}, scénario d'émission et
projection~; le langage calibré~: \emph{probable} de 66 à \qty{100}{\percent}, \emph{très probable} de 90 à
\qty{100}{\percent}, un intervalle $[x \text{ à } y]$ étant un intervalle très probable. L'influence humaine a
réchauffé le climat de \qty{1.07}{\celsius} entre 1850--1900 et 2010--2019. \textbf{Acquis des volumes de
physique}~: le corps noir et la loi de Stefan--Boltzmann~; la dilatation thermique. \textbf{Acquis des volumes
de biologie}~: le cycle du carbone, la fraction restée dans l'atmosphère, les puits de carbone et le bilan
carbone de la biosphère. \textbf{Acquis du \cref{ch:b2:deep-time-climate}}~: le cycle lent du
carbone, qui ne reprend l'excès de dioxyde de carbone qu'en centaines de milliers d'années.
\end{recall}
% ledger: ipcc:erf-2xco2, ipcc:log-co2, ipcc:ecs, ipcc:likelihood, ipcc:attribution

\section{Forçage quantifié}\label{sec:b2:anthropogenic-climate-change:forcing}

\begin{definition}[Forçage radiatif effectif]\label{def:b2:anthropogenic-climate-change:erf}
Le \emph{forçage radiatif effectif}\index{forçage radiatif effectif} (ERF) d'une perturbation
--- un changement de gaz à effet de serre, d'aérosols, du Soleil --- est la variation qu'elle
provoque dans le flux net descendant de rayonnement au sommet de l'atmosphère, en
\unit{W.m^{-2}}, une fois que l'atmosphère, les nuages et la surface continentale se sont
ajustés, mais avant que la température de surface globale ait réagi. Les ajustements retenus
sont ceux que cause l'agent lui-même, indépendamment de la température de surface.
\end{definition}
% ledger: ar6ch7:erf-def

La définition trace la frontière entre forçage et réponse selon ce qui dépend de la température de surface,
et non selon la rapidité. Le forçage instantané du dioxyde de carbone est la variation du flux juste après
l'ajout du gaz~; la stratosphère, qui se refroidit en quelques mois, le modifie une première fois~; la
troposphère, sa vapeur d'eau et ses nuages le modifient une seconde fois. Pour un doublement du dioxyde de
carbone, l'évaluation trouve \qty{3.75}{W.m^{-2}} une fois la stratosphère ajustée, et des ajustements dans
la troposphère qui se compensent en partie~: \qty{-0.60}{W.m^{-2}} dus à son réchauffement, +0.22 pour la
vapeur d'eau, +0.45 pour les nuages et +0.11 pour la surface continentale, soit \qty{+0.18}{W.m^{-2}} au
total. L'ERF vaut $3.75 + 0.18 = \qty{3.93}{W.m^{-2}}$, soit \qty{6}{\percent} de plus que les
\qty{3.71}{W.m^{-2}} de l'évaluation précédente (\cref{exo:b2:anthropogenic-climate-change:1}).
% ledger: ar6ch7:2xco2, ar6ch7:erf-def

Le forçage de chaque agent depuis 1750 a été reconstitué année par année. La figure ci-dessous le montre~:
le dioxyde de carbone à \qty{2.16}{W.m^{-2}} en 2019, les autres gaz à effet de serre --- méthane, protoxyde
d'azote, gaz halogénés et ozone --- à \qty{1.69}{W.m^{-2}} ensemble, et les aérosols à
\qty{-1.06}{W.m^{-2}}. Le forçage humain total est de \qty{2.72}{W.m^{-2}}, et son intervalle très probable,
de 1.96 à \qty{3.48}{W.m^{-2}}, est large surtout à cause des aérosols. Leur ERF sur 1750--2014 est évalué à
$-1.3$ [$-2.0$ à $-0.6$] \unit{W.m^{-2}} (\emph{confiance moyenne}), surtout par leur effet sur les nuages,
qui réfléchissent davantage la lumière lorsque les aérosols fournissent davantage de gouttelettes. Les
éruptions volcaniques ajoutent de brèves pointes négatives (\cref{ch:b2:eruptive-styles})~: la plus profonde
de l'enregistrement, \qty{-4.8}{W.m^{-2}} en 1816, suit l'éruption du Tambora en 1815.
% ledger: data:b4-ar6-erf, ipcc:erf-total, ipcc:erf-agents, ar6ch7:aerosol, usgs:msh-comparisons

\begin{omfigure}
\centering
\begin{tikzpicture}
\begin{axis}[omrecord, width=0.92\linewidth, height=6.2cm, xmin=1750, xmax=2020, ymin=-5.2,
  ymax=3.8, xlabel={année}, ylabel={ERF depuis 1750 (\unit{W.m^{-2}})},
  xtick={1750,1800,...,2000}, /pgf/number format/1000 sep={}, xtick pos=bottom, ytick pos=left,
  tick label style={font=\scriptsize}, label style={font=\small},
  legend style={at={(0.5,-0.24)}, anchor=north, font=\scriptsize, draw=black!30,
  legend columns=3, /tikz/every even column/.append style={column sep=6pt}},
  legend cell align=left]
% computed in figdata/bachelor-2/anthropogenic.py part erf; ledger: data:b4-ar6-erf
\addplot[name path=hi, draw=none, forget plot] table[x=year, y=hi] {figdata/out/bachelor-2-anthropogenic-erf.dat};
\addplot[name path=lo, draw=none, forget plot] table[x=year, y=lo] {figdata/out/bachelor-2-anthropogenic-erf.dat};
\addplot[black!15, forget plot] fill between[of=hi and lo];
\addplot[black!45, line width=0.4pt] table[x=year, y=natural] {figdata/out/bachelor-2-anthropogenic-erf.dat};
\addlegendentry{naturel (volcans, Soleil)}
\addplot[omThm, line width=1pt] table[x=year, y=co2] {figdata/out/bachelor-2-anthropogenic-erf.dat};
\addlegendentry{dioxyde de carbone}
\addplot[omThm!50!orange, line width=1pt, densely dashed] table[x=year, y=other_ghg] {figdata/out/bachelor-2-anthropogenic-erf.dat};
\addlegendentry{autres gaz à effet de serre et ozone}
\addplot[omDef, line width=1pt] table[x=year, y=aerosol] {figdata/out/bachelor-2-anthropogenic-erf.dat};
\addlegendentry{aérosols}
\addplot[black, line width=1.4pt] table[x=year, y=anthro] {figdata/out/bachelor-2-anthropogenic-erf.dat};
\addlegendentry{total humain (gris~: 5--95\,\%)}
\draw[black!40] (axis cs:1750,0) -- (axis cs:2020,0);
\node[anchor=west, font=\scriptsize, text=black!60, inner sep=1pt] at (axis cs:1817.5,-4.75) {$\leftarrow$ Tambora};
\end{axis}
\end{tikzpicture}
\par\smallskip{\small \omterm{def:b2:anthropogenic-climate-change:erf}{Forçage radiatif effectif} de 1750 à 2019, par rapport à 1750, par facteur. Le forçage
humain total atteint \qty{2.72}{W.m^{-2}} en 2019, son intervalle très probable étant ombré. Données~:
GIEC, sixième rapport d'évaluation, groupe de travail~I, données étendues de l'annexe~III (CC BY 4.0).}
\end{omfigure}
% values: anthro 2.7204 (1.9569-3.4819), co2 2.1563, other_ghg 1.6850, aerosol -1.0583, volcanic -4.8162 in 1816 (anthropogenic.py numbers)

\section{Bilan énergétique}\label{sec:b2:anthropogenic-climate-change:energy}

Un forçage positif fait gagner de l'énergie à la planète. En se réchauffant, elle rayonne davantage
vers l'espace, et le gain diminue~; le climat ne se stabilise que lorsque les deux se compensent.

\begin{definition}[Déséquilibre énergétique de la Terre~; paramètre de rétroaction climatique]\label{def:b2:anthropogenic-climate-change:eei}
\emph{Déséquilibre énergétique de la Terre}\index{déséquilibre énergétique de la Terre} $N$ est le flux net
d'énergie entrant dans le système climatique au sommet de l'atmosphère, moyenné sur le globe et sur une année
ou plus, en \unit{W.m^{-2}}~: énergie solaire absorbée moins infrarouge sortant. C'est le rythme auquel
l'océan, les continents, la glace et l'air gagnent de la chaleur, par mètre carré de la surface de la Terre.
\end{definition}

\begin{definition}[Paramètre de rétroaction climatique]\label{def:b2:anthropogenic-climate-change:lambda}
Le \emph{paramètre de rétroaction climatique}\index{paramètre de rétroaction climatique} $\lambda$ est la
variation du flux net descendant au sommet de l'atmosphère par degré de réchauffement global de la
surface, en \unit{W.m^{-2}.\celsius^{-1}}, causée par la réponse du système climatique à ce réchauffement.
Un climat stable a $\lambda < 0$~: le réchauffement accroît la perte d'énergie vers l'espace.
\end{definition}

Le volume scolaire écrivait le bilan avec la grandeur de $\lambda$, un nombre positif~;
ici le signe est conservé, comme dans les évaluations, qui écrivent le paramètre $\alpha$.
% ledger: ar6ch7:energy-budget

\begin{proposition}[Bilan énergétique global]\label{prop:b2:anthropogenic-climate-change:balance}
Soient $\Delta F$ le \omterm{def:b2:anthropogenic-climate-change:erf}{forçage radiatif effectif} et $\Delta T$ le réchauffement global de la surface
par rapport à un état d'équilibre. Si la réponse du flux sortant est linéaire en $\Delta T$,
\[
  \Delta N = \Delta F + \lambda\,\Delta T .
\]
À un nouvel équilibre $\Delta N = 0$, donc $\Delta T_{\text{éq}} = -\Delta F/\lambda$.
\end{proposition}
\begin{proof}
Par définition, le forçage modifie le flux net descendant de $\Delta F$ à température
de surface fixée, et la réponse au réchauffement le modifie de $\lambda\,\Delta T$ au
premier ordre~; le déséquilibre est leur somme. Le poser égal à zéro donne l'équilibre.
\end{proof}

Mesurer $N$ revient à mesurer la chaleur que gagne le système, et l'océan en retient presque toute la
part. De 1971 à 2018, le système climatique a gagné \qty{435}{ZJ} (zettajoules, \qty{e21}{J})~: l'océan
\qty{91}{\percent}, soit \qty{242}{ZJ} dans ses \qty{700}{m} supérieurs et \qty{154}{ZJ} au-dessous~;
les terres \qty{5}{\percent}, la glace qui fond \qty{3}{\percent}, l'atmosphère \qty{1}{\percent}. Le
déséquilibre moyen était de \qty{0.57}{W.m^{-2}} sur 1971--2018, et il a augmenté~:
\qty{0.79}{W.m^{-2}} [0.52 à 1.06] sur 2006--2018 (\cref{exo:b2:anthropogenic-climate-change:2}).
% ledger: ar6ch7:inventory, ipcc:heating

\begin{omfigure}
\centering
\begin{tikzpicture}[font=\small, >=Stealth]
\draw[fill=omDef!8, draw=omDef!60] (-3.2,0) rectangle (3.2,1.6);
\node[align=center] at (0,1.05) {système climatique};
\node[align=center, font=\scriptsize] at (0,0.5) {océan \qty{91}{\percent}, terres \qty{5}{\percent}, glace \qty{3}{\percent}, atmosphère \qty{1}{\percent}};
\draw[->, line width=2.2pt, omThm] (-2.0,3.3) -- node[pos=0.12, left=3pt, align=right, font=\scriptsize] {forçage $\Delta F$} (-2.0,1.7);
\draw[->, line width=2.2pt, omDef] (2.0,1.7) -- node[pos=0.88, right=3pt, align=left, font=\scriptsize] {réponse $\lambda\,\Delta T$} (2.0,3.3);
\draw[->, line width=1.4pt, black!70] (0,-0.1) -- node[right=3pt, align=left, font=\scriptsize] {stocké~: $\Delta N = \Delta F + \lambda\,\Delta T$} (0,-1.0);
\draw[black!40, dashed] (-3.4,2.4) -- (3.4,2.4) node[right, font=\scriptsize, text=black!60] {sommet de l'atmosphère};
\end{tikzpicture}
\par\smallskip{\small Le bilan énergétique d'un climat soumis à un forçage (schéma).
Le forçage apporte de l'énergie au sommet de l'atmosphère, la réponse de la planète
plus chaude en retire une partie ($\lambda\,\Delta T$ est négatif), et la différence
est stockée, surtout dans l'océan. Parts de la chaleur gagnée de 1971 à 2018.}
\end{omfigure}
% ledger: ar6ch7:inventory

\begin{method}[Une sensibilité climatique à partir du bilan énergétique historique]\label{met:b2:anthropogenic-climate-change:energy-budget}
\begin{enumerate}
  \item Choisir deux périodes, une de référence proche de l'équilibre et une récente~; prendre la variation du
        forçage $\Delta F$, du déséquilibre $\Delta N$ et de la température globale $\Delta T$.
  \item Résoudre le bilan pour le paramètre de rétroaction~: $\lambda = (\Delta N - \Delta F)/\Delta T$.
  \item La sensibilité à un doublement vaut alors $-F_{2\times}/\lambda$~; le réchauffement au moment du
        forçage, ramené à $F_{2\times}$, vaut $F_{2\times}\Delta T/\Delta F$.
\end{enumerate}
\end{method}

De 1850--1900 à 2006--2019, l'évaluation prend $\Delta N = \qty{0.59}{W.m^{-2}}$, $\Delta T =
\qty{1.03}{\celsius}$ et $\Delta F = \qty{2.20}{W.m^{-2}}$. Alors $\lambda = (0.59 - 2.20)/1.03 =
\qty{-1.56}{W.m^{-2}.\celsius^{-1}}$, et la sensibilité vaut $3.93/1.56 = \qty{2.5}{\celsius}$, la valeur
retenue par l'évaluation, avec un intervalle très probable de 1.6 à \qty{4.8}{\celsius}. Mais il s'agit
d'une sensibilité \emph{effective}, pour les rétroactions du réchauffement jusqu'ici. Le réchauffement a
été le plus fort là où les rétroactions sont les plus négatives, dans le Pacifique tropical occidental,
et plus faible sur le Pacifique oriental et l'océan Austral~; à mesure que ce motif évolue, le paramètre
de rétroaction devrait devenir moins négatif (\emph{confiance élevée}), de sorte que la sensibilité à
long terme est supérieure à l'estimation historique (\cref{exo:b2:anthropogenic-climate-change:7}).
% ledger: ar6ch7:energy-budget, ar6ch7:pattern, ipcc:erf-2xco2
% values: lambda_energy -1.563, ecs_energy 2.514 (anthropogenic.py numbers)

\section{Rétroactions et sensibilité}\label{sec:b2:anthropogenic-climate-change:feedbacks}

\begin{definition}[Réponse de Planck~; rétroaction de la vapeur d'eau, rétroaction du gradient thermique, rétroaction des nuages]\label{def:b2:anthropogenic-climate-change:feedbacks}
Le paramètre de rétroaction se décompose en contributions des grandeurs qui varient avec la température de
surface. La \emph{réponse de Planck}\index{réponse de Planck} est le rayonnement infrarouge supplémentaire
émis vers l'espace lorsque la surface et toute la troposphère se réchauffent de la même quantité. La
\emph{rétroaction de la vapeur d'eau}\index{rétroaction de la vapeur d'eau} est la variation du flux due à
l'augmentation de la vapeur d'eau atmosphérique avec la température. La \emph{rétroaction du gradient
thermique}\index{rétroaction du gradient thermique} est la variation du flux due à un réchauffement non
uniforme selon l'altitude. La \emph{rétroaction des nuages}\index{rétroaction des nuages} est la variation
du flux due à la variation de la quantité, de l'altitude et de la brillance des nuages.
\end{definition}

\begin{proposition}[Réponse de Planck]\label{prop:b2:anthropogenic-climate-change:planck}
Si la Terre émet vers l'espace une fraction fixe $1 - g$ du rayonnement infrarouge
$\sigma T^4$ de sa surface, à la température $T$, sa \omterm{def:b2:anthropogenic-climate-change:feedbacks}{réponse de Planck} est
\[
  \lambda_{\text P} = -(1 - g)\,4\sigma T^3 .
\]
Avec $g = 0.4$ et $T = \qty{288}{K}$, $\lambda_{\text P} = \qty{-3.25}{W.m^{-2}.K^{-1}}$.
\end{proposition}
\begin{proof}
Le flux sortant est $(1 - g)\sigma T^4$~; sa dérivée par rapport à $T$ est $(1 - g)\,4\sigma
T^3$, et une perte supplémentaire est une contribution négative au flux net descendant. Numériquement $0.6
\times 4 \times \num{5.670e-8} \times 288^3 = 3.25$.
\end{proof}
% ledger: ar6ch7:planck, const:sigma, ipcc:budget-surface
% values: planck -3.2509 (anthropogenic.py numbers)

La fraction $g$ est l'effet de serre normalisé~: l'atmosphère retient \qty{159}{W.m^{-2}} des \qty{398}{W.m^{-2}}
que la surface émet, soit $g \approx 0.4$. L'évaluation, avec des codes de rayonnement détaillés, trouve $-3.22$
[$-3.4$ à $-3.0$] (\emph{confiance élevée}). La vapeur d'eau suit la loi de Clausius--Clapeyron à l'échelle
globale, et davantage de vapeur absorbe davantage d'infrarouge~: la \omterm{def:b2:anthropogenic-climate-change:feedbacks}{rétroaction de la vapeur d'eau} est fortement
positive, \qty{1.85}{W.m^{-2}.\celsius^{-1}} d'après les satellites. Sous les tropiques, la haute troposphère se
réchauffe plus que la surface, qui rayonne alors davantage~: la \omterm{def:b2:anthropogenic-climate-change:feedbacks}{rétroaction du gradient thermique} y est négative,
positive près des pôles, et d'environ $-0.50$ à l'échelle globale. Les deux sont évaluées ensemble, \num{1.30}
[1.1 à 1.5], car ce qu'un modèle met dans l'une tend à le prendre à l'autre. La neige et la glace qui fondent
assombrissent la surface, +0.35~; les nuages ajoutent \num{0.42} [$-0.10$ à $0.94$], la plus grande incertitude.
Une \omterm{def:b2:anthropogenic-climate-change:feedbacks}{rétroaction des nuages} nettement négative est \emph{très improbable}.
% ledger: ar6ch7:planck, ar6ch7:wv-lr, ar6ch7:table710, ipcc:cloud-negative, ipcc:budget-surface

\begin{omfigure}
\centering
\begin{tikzpicture}
\begin{axis}[xbar, width=0.78\linewidth, height=5.4cm, xmin=-3.7, xmax=1.8, ymin=0.4, ymax=6.6,
  ytick={1,...,6}, yticklabels={net, biogéophysique et chimique, nuages, albédo de surface,
  vapeur d'eau + gradient thermique, Planck}, y tick label style={font=\scriptsize},
  tick label style={font=\scriptsize}, xlabel={paramètre de rétroaction (\unit{W.m^{-2}.\celsius^{-1}})},
  label style={font=\small}, bar width=7pt, axis x line=bottom, axis y line=left,
  axis line style={-}, xmajorgrids, grid style={gray!22, line width=0.3pt}, extra x ticks={0},
  extra x tick style={grid style={black!60, line width=0.5pt}}]
% ledger: ar6ch7:table710 (central values and very likely ranges)
\addplot[fill=omDef!55, draw=omDef!80!black, error bars/.cd, x dir=both, x explicit,
  error bar style={black!70}] coordinates {
  (-3.22,6) += (0.22,0) -= (0.18,0)
  (-0.01,2) += (0.26,0) -= (0.26,0)
  (-1.16,1) += (0.65,0) -= (0.65,0)};
\addplot[fill=omThm!55, draw=omThm!80!black, error bars/.cd, x dir=both, x explicit,
  error bar style={black!70}] coordinates {
  (1.30,5) += (0.20,0) -= (0.20,0)
  (0.35,4) += (0.25,0) -= (0.25,0)
  (0.42,3) += (0.52,0) -= (0.52,0)};
\end{axis}
\end{tikzpicture}
\par\smallskip{\small Les rétroactions climatiques évaluées par le GIEC (sixième rapport d'évaluation, groupe
de travail~I, tableau 7.10), valeurs centrales et intervalles très probables. Les rétroactions positives
amplifient un réchauffement (rouge), les négatives l'amortissent (bleu)~; la rétroaction nette est leur somme.}
\end{omfigure}

\begin{definition}[Sensibilité climatique à l'équilibre]\label{def:b2:anthropogenic-climate-change:ecs}
La \emph{sensibilité climatique à l'équilibre}\index{sensibilité climatique à l'équilibre} (ECS) est le
réchauffement global de la surface, au nouvel équilibre, causé par un doublement du dioxyde de carbone
de l'air par rapport à son niveau préindustriel, les calottes glaciaires étant maintenues fixes.
\end{definition}

\begin{proposition}[Sensibilité et rétroaction]\label{prop:b2:anthropogenic-climate-change:ecs}
Si le paramètre de rétroaction est constant, $\text{ECS} = -F_{2\times}/\lambda$. Avec $F_{2\times} =
\qty{3.93}{W.m^{-2}}$ et le $\lambda = \qty{-1.16}{W.m^{-2}.\celsius^{-1}}$ net des rétroactions
ci-dessus, $\text{ECS} = \qty{3.4}{\celsius}$~; la seule \omterm{def:b2:anthropogenic-climate-change:feedbacks}{réponse de Planck} donnerait $3.93/3.22 =
\qty{1.2}{\celsius}$.
\end{proposition}
\begin{proof}
Appliquer la \cref{prop:b2:anthropogenic-climate-change:balance} à l'équilibre avec $\Delta F =
F_{2\times}$.
\end{proof}
% ledger: ar6ch7:table710, ipcc:erf-2xco2
% values: ecs_process 3.388 (anthropogenic.py numbers)

Les rétroactions ne sont qu'un élément de preuve~; le réchauffement depuis 1850, le bilan énergétique
et les climats du passé en sont d'autres. Combinés, ils donnent l'ECS évaluée à \qty{3}{\celsius},
\emph{probable} de 2.5 à \qty{4}{\celsius} (\emph{confiance élevée}), \emph{très probable} de 2 à
\qty{5}{\celsius}. Comme l'ECS varie en $1/\lambda$, une incertitude symétrique sur $\lambda$ donne
un intervalle d'ECS étiré vers les valeurs élevées (\cref{exo:b2:anthropogenic-climate-change:12}).
% ledger: ipcc:ecs

\begin{definition}[Réponse climatique transitoire]\label{def:b2:anthropogenic-climate-change:tcr}
La \emph{réponse climatique transitoire}\index{réponse climatique transitoire} (TCR) est le réchauffement
global de la surface au moment où le dioxyde de carbone de l'air a doublé, lorsqu'il augmente de
\qty{1}{\percent} par an à partir de son niveau préindustriel, soit environ soixante-dix ans après le début.
\end{definition}

La TCR est plus petite que l'ECS, car l'océan absorbe encore de la chaleur au moment où le doublement est
atteint. Sa meilleure estimation est \qty{1.8}{\celsius}, \emph{probable} de 1.4 à \qty{2.2}{\celsius}
(\emph{confiance élevée})~; le seul bilan énergétique historique donne aussi $3.93 \times 1.03/2.20 =
\qty{1.8}{\celsius}$ également.
% ledger: ar6ch7:tcr, ar6ch7:energy-budget
% values: tcr_energy 1.840 (anthropogenic.py numbers)

\section{Détection et attribution}\label{sec:b2:anthropogenic-climate-change:attribution}

\begin{definition}[Détection~; attribution~; empreinte climatique]\label{def:b2:anthropogenic-climate-change:da}
\emph{Détection}\index{détection} d'un changement climatique consiste à montrer que le climat a changé
au sens statistique défini, au-delà de ce que sa variabilité interne peut produire, sans en donner la
cause. \emph{Attribution}\index{attribution} consiste à évaluer les contributions de plusieurs causes à
un changement détecté, avec une vraisemblance attribuée. Une \emph{empreinte climatique}\index{empreinte climatique}
est le motif, dans l'espace et dans le temps, de la réponse à un facteur tel que simulé par
les modèles climatiques, grâce auquel ce facteur est reconnu dans les observations.
\end{definition}

Des facteurs différents laissent des empreintes différentes. Davantage de gaz à effet de serre réchauffe la troposphère
et refroidit la stratosphère, car une couverture infrarouge plus épaisse fait rayonner davantage les couches supérieures
vers l'espace~; un Soleil plus brillant réchaufferait les deux. Les observations montrent le premier motif~: la
troposphère s'est réchauffée depuis au moins les années 1950, et il est \emph{virtuellement certain} que la stratosphère
s'est refroidie. Il est \emph{très probable} que les gaz à effet de serre d'origine humaine ont été le principal moteur
du réchauffement troposphérique depuis 1979, et \emph{extrêmement probable} que l'appauvrissement de la couche d'ozone
ait été le principal moteur du refroidissement de la basse stratosphère de 1979 au milieu des années 1990.
% ledger: ar6ch3:upper-air

\begin{method}[Empreinte optimale]\label{met:b2:anthropogenic-climate-change:fingerprinting}
\begin{enumerate}
  \item Simuler la réponse à chaque facteur séparément~: les empreintes $x_1, \dots, x_k$, sous forme
        de vecteurs de valeurs dans l'espace et dans le temps.
  \item Écrire les observations sous la forme $y = \sum_i \beta_i x_i + \varepsilon$, où $\varepsilon$ est la
        variabilité interne, dont les statistiques sont tirées de longues simulations sans forçage.
  \item Estimer les facteurs d'échelle $\beta_i$ par moindres carrés, pondérés par l'inverse de la
        variabilité, pour que les composantes les plus bruitées comptent le moins, avec leurs intervalles de confiance.
  \item Un $\beta_i$ nettement supérieur à zéro signifie que la réponse au facteur $i$ est détectée~;
        un intervalle qui contient 1 et exclut 0 signifie que la réponse simulée est compatible avec les
        observations, et que le changement peut être attribué en partie à ce facteur.
\end{enumerate}
\end{method}
% ledger: ar6ch3:fingerprint

Appliquée à la température de surface, la méthode donne l'\omterm{def:b2:anthropogenic-climate-change:da}{attribution} rappelée au début~: sur les
\qty{1.06}{\celsius} observés entre 1850--1900 et 2010--2019, l'influence humaine a \emph{probablement}
causé de 0.8 à \qty{1.3}{\celsius}, meilleure estimation \qty{1.07}{\celsius}~; les facteurs naturels
et la variabilité interne ne contribuent pas plus de quelques dixièmes de degré dans un sens ou dans
l'autre (\cref{exo:b2:anthropogenic-climate-change:9,exo:b2:anthropogenic-climate-change:10}).
% ledger: ipcc:attribution, ipcc:observed-2010-2019

\section{Scénarios et émissions cumulées}\label{sec:b2:anthropogenic-climate-change:scenarios}

\begin{definition}[Trajectoire socio-économique partagée]\label{def:b2:anthropogenic-climate-change:ssp}
Une \emph{trajectoire socio-économique partagée}\index{trajectoire socio-économique partagée} (SSP) est
l'un de cinq récits sur la manière dont la population, les économies et la technologie pourraient évoluer
au cours du siècle. Les scénarios d'émission de l'évaluation sont notés SSP$x$-$y$~: $x$ désigne la
trajectoire, et $y$ le forçage radiatif approximatif, en \unit{W.m^{-2}}, atteint par le scénario en 2100.
\end{definition}
% ledger: ar6spm:ssp

Les cinq trajectoires principales vont de SSP1-1.9 à SSP5-8.5. L'évaluation de la science physique
les prend telles quelles~: elle ne pondère ni leur vraisemblance ni les récits eux-mêmes. Pour
chacune, les modèles climatiques donnent un réchauffement projeté pour 2081--2100 par rapport à
1850--1900, de \qty{1.4}{\celsius} (très probable de 1.0 à 1.8) pour SSP1-1.9 à \qty{4.4}{\celsius}
(de 3.3 à 5.7) pour SSP5-8.5. Les projections obéissent à une loi simple commune.
% ledger: ar6spm:ssp, ipcc:proj-temp

\begin{definition}[Réponse climatique transitoire aux émissions cumulées de CO$_2$]\label{def:b2:anthropogenic-climate-change:tcre}
La \emph{réponse climatique transitoire aux émissions cumulées de
\ce{CO2}}\index{réponse climatique transitoire aux émissions cumulées de CO$_2$} (TCRE)
est le réchauffement global de la surface par unité de dioxyde de carbone émis au total
depuis la période préindustrielle, en \unit{\celsius} par \qty{1000}{Gt} de \ce{CO2}.
\end{definition}

\begin{proposition}[Le réchauffement suit les émissions cumulées]\label{prop:b2:anthropogenic-climate-change:tcre}
À une bonne approximation, le réchauffement causé par le dioxyde de carbone est
proportionnel au total émis, quelle que soit la trajectoire des émissions~:
$\Delta T \approx \text{TCRE} \times E$, avec $E$ les émissions cumulées.
\end{proposition}
\admitted
% ledger: ar6spm:tcre, ar6ch5:tcre-linear

La loi n'est pas évidente. Le forçage ne croît que comme le logarithme de la concentration, si bien que chaque tonne
ajoutée réchauffe moins~; mais l'océan absorbe la chaleur et le carbone de moins en moins efficacement à mesure que
les émissions s'accumulent, si bien qu'une plus grande part de chaque tonne reste dans l'air et que son réchauffement
est moins retardé. Les deux effets se compensent presque, et toutes les études s'accordent à en faire l'origine de la
loi. La relation est presque linéaire (\emph{confiance élevée}), et chaque \qty{1000}{Gt} de dioxyde de carbone
\emph{cause probablement} un réchauffement de 0.27 à \qty{0.63}{\celsius}, meilleure estimation \qty{0.45}{\celsius},
soit \qty{1.65}{\celsius} par \qty{1000}{PgC} de carbone (\cref{exo:b2:anthropogenic-climate-change:4}). De 1850 à
2019, $\num{2390} \pm \qty{240}{Gt}$ de dioxyde de carbone ont été émis.
% ledger: ar6ch5:tcre-linear, ar6spm:tcre, ar6spm:budget

\begin{definition}[Budget carbone restant]\label{def:b2:anthropogenic-climate-change:budget}
Le \emph{budget carbone restant}\index{budget carbone restant} pour une limite de
température est le dioxyde de carbone cumulé qui peut encore être émis, à partir d'une
date donnée, le réchauffement global restant inférieur à cette limite avec une probabilité
donnée, le réchauffement causé par les autres facteurs humains étant pris en compte.
\end{definition}

\begin{method}[Le budget carbone restant à partir de la TCRE]\label{met:b2:anthropogenic-climate-change:budget}
\begin{enumerate}
  \item Retrancher de la limite le réchauffement humain déjà atteint, \qty{1.07}{\celsius} pour
        2010--2019.
  \item Retrancher le réchauffement que les autres facteurs ajouteront d'ici à ce que les émissions de
        dioxyde de carbone atteignent le zéro net.
  \item Diviser par la TCRE~: c'est le budget à partir du milieu de la période 2010--2019.
  \item Retrancher les émissions depuis, et une marge pour les rétroactions lentes du système Terre que les
        modèles omettent (le dégel du pergélisol est le plus important).
\end{enumerate}
\end{method}

L'évaluation combine cinq composantes de cette manière~: la TCRE~; le réchauffement historique~; un
réchauffement supplémentaire dû aux autres facteurs, d'environ 0.1 à \qty{0.2}{\celsius}~; le réchauffement
qui peut se poursuivre après l'arrêt des émissions, évalué à zéro en moyenne (\emph{probable} dans $\pm
\qty{0.19}{\celsius}$)~; et les rétroactions du système Terre, d'environ \qty{26}{Gt} de dioxyde de carbone
par degré. Pour \qty{1.5}{\celsius}, soit \qty{0.43}{\celsius} au-dessus de 2010--2019, le budget restant à
partir du début de 2020 est de \qty{500}{Gt} de dioxyde de carbone avec une probabilité de
\qty{50}{\percent}, et \qty{1350}{Gt} pour \qty{2}{\celsius}. Les choix faits sur les autres facteurs
déplacent ces valeurs d'au moins \qty{220}{Gt}~; le problème du week-end en reconstruit la première.
% ledger: ar6ch5:budget-parts, ar6spm:budget

\begin{dated}{2025}
% ledger: data:co2-mlo, data:gistemp (values read from the datasets as fetched 2026-10-22; offset computed)
La moyenne annuelle du dioxyde de carbone au Mauna Loa a atteint \qty{427.35}{ppm} en 2025, soit un
forçage de $3.93\log_2(427.35/278) = \qty{2.44}{W.m^{-2}}$ selon la loi logarithmique. Dans la série de
températures de la NASA, 2024 est l'année la plus chaude, \qty{1.50}{\celsius} au-dessus de la moyenne
de 1880--1900, les premières années de cette série~; 2025 s'est établie à \qty{1.40}{\celsius}.
\end{dated}

\begin{omfigure}
\centering
\begin{tikzpicture}
\begin{axis}[omrecord, width=0.86\linewidth, height=5.8cm, xmin=0, xmax=4000, ymin=0, ymax=2.6,
  xlabel={émissions cumulées de \ce{CO2} depuis 1850 (\unit{Gt})},
  ylabel={réchauffement depuis 1850--1900 (\unit{\celsius})}, xtick={0,500,...,4000},
  /pgf/number format/1000 sep={\,}, xtick pos=bottom, ytick pos=left, clip=false,
  tick label style={font=\scriptsize}, label style={font=\small}]
% computed in figdata/bachelor-2/anthropogenic.py part tcre; ledger: ar6spm:tcre, ar6spm:budget
\addplot[name path=hi, draw=none, forget plot] table[x=E_GtCO2, y=dT_hi] {figdata/out/bachelor-2-anthropogenic-tcre.dat};
\addplot[name path=lo, draw=none, forget plot] table[x=E_GtCO2, y=dT_lo] {figdata/out/bachelor-2-anthropogenic-tcre.dat};
\addplot[omThm!18, forget plot] fill between[of=hi and lo];
\addplot[omThm, line width=1.3pt] table[x=E_GtCO2, y=dT_mid] {figdata/out/bachelor-2-anthropogenic-tcre.dat};
\draw[densely dashed, black!60] (axis cs:0,1.5) -- (axis cs:4000,1.5);
\draw[densely dashed, black!60] (axis cs:0,2.0) -- (axis cs:4000,2.0);
\node[anchor=north west, font=\scriptsize, text=black!70] at (axis cs:0,1.5) {\qty{1.5}{\celsius}};
\node[anchor=north west, font=\scriptsize, text=black!70] at (axis cs:0,2.0) {\qty{2}{\celsius}};
\fill[black] (axis cs:2390,1.07) circle (2.2pt);
\draw[black, line width=0.8pt] (axis cs:2150,1.07) -- (axis cs:2630,1.07);
\draw[black, line width=0.8pt] (axis cs:2390,0.8) -- (axis cs:2390,1.3);
\node[anchor=north west, font=\scriptsize, align=left] at (axis cs:2440,1.02) {1850--2019~:\\\qty{2390}{Gt}, \qty{1.07}{\celsius}};
\draw[<->, omDef, thick] (axis cs:2390,1.62) -- (axis cs:2890,1.62);
\node[anchor=south, font=\scriptsize, text=omDef, align=center] at (axis cs:2640,1.64) {AR6~: \qty{500}{Gt}};
\node[anchor=north west, font=\scriptsize, text=omThm, align=left] at (axis cs:150,2.5) {intervalle probable de la TCRE,\\0.27 à \qty{0.63}{\celsius} par \qty{1000}{Gt}};
\end{axis}
\end{tikzpicture}
\par\smallskip{\small Réchauffement en fonction des émissions cumulées de dioxyde de carbone, pour la
TCRE évaluée, \qty{0.45}{\celsius} par \qty{1000}{Gt} (trait) et son intervalle probable (zone ombrée),
avec le réchauffement humain et les émissions à ce jour (point, intervalles probables en barres). La
flèche marque le budget restant pour \qty{1.5}{\celsius} à \qty{50}{\percent} à partir de 2020~: plus
petit que ce que donnerait la seule lecture du trait, car les autres facteurs réchauffent aussi.}
\end{omfigure}


\section{Niveau de la mer}\label{sec:b2:anthropogenic-climate-change:sea-level}

\begin{definition}[Bilan du niveau marin~; élévation thermostérique du niveau marin]\label{def:b2:anthropogenic-climate-change:sea-level}
Le \emph{bilan du niveau marin}\index{bilan du niveau marin} est la comparaison de la variation observée du
niveau moyen global de la mer avec la somme de ses contributions estimées~: l'expansion de l'océan qui se
réchauffe, la fonte des glaciers et des calottes du Groenland et de l'Antarctique, et la variation de l'eau
stockée sur les continents. L'\emph{élévation thermostérique du niveau marin}\index{élévation thermostérique du niveau marin}
est la part due à la dilatation thermique de l'eau de mer, à masse constante.
\end{definition}

\begin{proposition}[Dilatation thermique d'une couche]\label{prop:b2:anthropogenic-climate-change:thermosteric}
Une couche d'océan d'épaisseur $H$, de coefficient de dilatation thermique $\alpha$, qui se réchauffe
uniformément de $\Delta T$ à masse et surface constantes, s'élève de $\Delta h = \alpha\,\Delta T\,H$.
\end{proposition}
\begin{proof}
La masse $\rho H A$ d'une colonne d'aire $A$ est constante, donc $\Delta H/H = -\Delta\rho/\rho =
\alpha\,\Delta T$ au premier ordre, avec $\alpha = -(1/\rho)\,\partial\rho/\partial T$.
\end{proof}
% ledger: teos10:alpha-beta

Dans l'océan réel, $\alpha$ croît avec la température et la pression, et le réchauffement est concentré
près de la surface. L'évaluation résume le résultat par une efficacité d'expansion de \qty{0.113}{m}
par \qty{e24}{J} de chaleur ajoutée à l'océan, soit \qty{0.617}{m} par degré de réchauffement de tout
l'océan (\emph{confiance moyenne}). Sur sa profondeur moyenne de \qty{3682}{m}, cela donne un $\alpha$
effectif de $0.617/3682 = \qty{1.68e-4}{\celsius^{-1}}$ (\cref{exo:b2:anthropogenic-climate-change:8}).
% ledger: ar6ch9:expansion, charette:ocean-volume

Le bilan, représenté dans la figure ci-dessous, se ferme dans les limites de ses incertitudes. De 1971 à 2018, les
contributions s'additionnent à \qty{2.00}{mm.an^{-1}}, contre une hausse observée de \qty{2.33}{mm.an^{-1}}~; de 2006 à 2018,
elles s'additionnent à 3.61 contre 3.69. La dilatation thermique a été le plus grand terme isolé, la moitié de la somme sur
1971--2018, mais la glace rattrape son retard~: de 2006 à 2018, les glaciers et les calottes glaciaires ont donné ensemble
\qty{1.62}{mm.an^{-1}}, davantage que la dilatation. L'océan continue aussi de se réchauffer en profondeur longtemps après
que la surface s'est stabilisée, et les calottes glaciaires réagissent sur plusieurs siècles~; le niveau de la mer est engagé
à monter pendant des siècles à des millénaires et restera élevé pendant des milliers d'années (\emph{confiance élevée}). Sur
les deux mille prochaines années, il devrait monter d'environ 2 à \qty{3}{m} si le réchauffement reste à \qty{1.5}{\celsius},
de 2 à \qty{6}{m} à \qty{2}{\celsius} et de 19 à \qty{22}{m} à \qty{5}{\celsius} (\emph{confiance faible}).
% ledger: ar6ch9:table95, ar6spm:sea-commit

\begin{omfigure}
\centering
\begin{tikzpicture}
\begin{axis}[ybar stacked, width=0.72\linewidth, height=6.2cm, ymin=0, ymax=4.3, bar width=26pt,
  xtick=data, symbolic x coords={1971--2018, 1993--2018, 2006--2018}, xticklabels={1971--2018, 1993--2018, 2006--2018}, enlarge x limits=0.25,
  ylabel={vitesse de la hausse (\unit{mm.an^{-1}})}, tick label style={font=\scriptsize},
  label style={font=\small}, ytick pos=left, xtick pos=bottom,
  legend style={at={(1.03,0.5)}, anchor=west, font=\scriptsize, draw=black!30},
  legend cell align=left, reverse legend]
% ledger: ar6ch9:table95
\addplot[fill=omDef!70, draw=omDef!90!black] coordinates {(1971--2018,1.01) (1993--2018,1.31) (2006--2018,1.39)};
\addlegendentry{dilatation thermique}
\addplot[fill=omDef!30, draw=omDef!70!black] coordinates {(1971--2018,0.44) (1993--2018,0.55) (2006--2018,0.62)};
\addlegendentry{glaciers}
\addplot[fill=omThm!55, draw=omThm!80!black] coordinates {(1971--2018,0.25) (1993--2018,0.43) (2006--2018,0.63)};
\addlegendentry{Groenland}
\addplot[fill=omThm!25, draw=omThm!70!black] coordinates {(1971--2018,0.14) (1993--2018,0.25) (2006--2018,0.37)};
\addlegendentry{Antarctique}
\addplot[fill=black!20, draw=black!50] coordinates {(1971--2018,0.15) (1993--2018,0.31) (2006--2018,0.60)};
\addlegendentry{eau terrestre}
\node[fill=black, rotate=45, minimum size=5pt, inner sep=0pt] at (axis cs:1971--2018,2.33) {};
\node[fill=black, rotate=45, minimum size=5pt, inner sep=0pt] at (axis cs:1993--2018,3.25) {};
\node[fill=black, rotate=45, minimum size=5pt, inner sep=0pt] at (axis cs:2006--2018,3.69) {};
\end{axis}
\end{tikzpicture}
\par\smallskip{\small Le \omterm{def:b2:anthropogenic-climate-change:sea-level}{bilan du niveau marin}~: contributions observées à la hausse du niveau
moyen global de la mer, sous forme de vitesses sur trois périodes se terminant en 2018, empilées,
et la vitesse observée (losanges). L'écart entre la pile et le losange est la part que le bilan
ne ferme pas. Valeurs~: GIEC, sixième rapport d'évaluation, groupe de travail~I, tableau 9.5.}
\end{omfigure}

\section*{Exercices}\label{sec:b2:anthropogenic-climate-change:exercises}

Données des exercices, sauf indication contraire~: forçage du dioxyde de carbone doublé $F_{2\times} =
\qty{3.93}{W.m^{-2}}$ à partir de $C_0 = \qty{278}{ppm}$~; constante de Stefan--Boltzmann $\sigma =
\qty{5.670e-8}{W.m^{-2}.K^{-4}}$~; rayon de la Terre \qty{6371}{km}~; une année \qty{3.156e7}{s}~;
masses molaires du carbone et du dioxyde de carbone 12.011 et \qty{44.009}{g.mol^{-1}}~; les
rétroactions et le \omterm{def:b2:anthropogenic-climate-change:sea-level}{bilan du niveau marin} du chapitre.
% ledger: ipcc:erf-2xco2, const:sigma, earth:Rmean, ciaaw:abridged-2024, ar6ch7:table710, ar6ch9:table95

\begin{exercise}[$\star$]\label{exo:b2:anthropogenic-climate-change:1}
Le forçage du dioxyde de carbone doublé est \qty{3.75}{W.m^{-2}} une fois la stratosphère
ajustée~; la troposphère s'ajuste alors de $-0.60$ (température), +0.22 (vapeur d'eau),
+0.45 (nuages) et +0.11 (surface continentale) \unit{W.m^{-2}}. (a) Trouver l'ERF. (b) De
quel pourcentage dépasse-t-il les \qty{3.71}{W.m^{-2}} de l'évaluation précédente~? (c) Quel
ajustement est négatif, et pourquoi une troposphère plus chaude réduit-elle le forçage~?
\end{exercise}
% ledger: ar6ch7:2xco2

\begin{exercise}[$\star$]\label{exo:b2:anthropogenic-climate-change:2}
Le système climatique a gagné \qty{434.9}{ZJ} de 1971 à 2018 et \qty{152.4}{ZJ} de 2006 à
2018, comptés entre les milieux de la première et de la dernière année. (a) Trouver la
surface de la Terre. (b) Convertir chaque gain en un déséquilibre moyen en \unit{W.m^{-2}}.
(c) L'océan a pris \qty{91}{\percent}~: combien de zettajoules sur 1971--2018~?
\end{exercise}
% ledger: ar6ch7:inventory, earth:Rmean

\begin{exercise}[$\star$]\label{exo:b2:anthropogenic-climate-change:3}
(a) Additionner les valeurs centrales des rétroactions de la figure, et comparer à la valeur
nette. (b) Trouver l'ECS avec la valeur nette, et le réchauffement avec la seule \omterm{def:b2:anthropogenic-climate-change:feedbacks}{réponse de
Planck}. (c) De quel facteur les autres rétroactions amplifient-elles le réchauffement de Planck~?
\end{exercise}

\begin{exercise}[$\star$]\label{exo:b2:anthropogenic-climate-change:4}
(a) Convertir une TCRE de \qty{0.45}{\celsius} par \qty{1000}{Gt} de dioxyde de carbone en
degrés par \qty{1000}{PgC} de carbone, et faire de même pour son intervalle probable, de
0.27 à 0.63. (b) Convertir les \qty{2390}{Gt} de dioxyde de carbone émis de 1850 à 2019 en
pétagrammes de carbone. (c) Quel réchauffement la TCRE donne-t-elle pour ces émissions~?
\end{exercise}
% ledger: ar6spm:tcre, ar6spm:budget

\begin{exercise}[$\star\star$]\label{exo:b2:anthropogenic-climate-change:5}
(a) La surface émet \qty{398}{W.m^{-2}}, et l'effet de serre en retient \qty{159}{W.m^{-2}}.
Trouver $g$ et la température de surface d'un corps noir émettant \qty{398}{W.m^{-2}}. (b)
Trouver la \omterm{def:b2:anthropogenic-climate-change:feedbacks}{réponse de Planck} à \qty{288}{K}, et à \qty{298}{K} avec le même $g$. (c) La
seule \omterm{def:b2:anthropogenic-climate-change:feedbacks}{réponse de Planck} rend-elle une Terre plus chaude plus ou moins sensible~?
\end{exercise}
% ledger: ipcc:budget-surface, ar6ch7:planck

\begin{exercise}[$\star\star$]\label{exo:b2:anthropogenic-climate-change:6}
Les satellites donnent une \omterm{def:b2:anthropogenic-climate-change:feedbacks}{rétroaction de la vapeur d'eau} de $1.85 \pm 0.32$ et les études une \omterm{def:b2:anthropogenic-climate-change:feedbacks}{rétroaction
du gradient thermique} de $-0.50 \pm 0.20$ \unit{W.m^{-2}.\celsius^{-1}} (un écart type). (a) Trouver leur
somme. (b) Si les deux étaient indépendantes, quel serait l'écart type de la somme~? (c) Les modèles
donnent à la somme un écart type de seulement \num{0.10}. Qu'est-ce que cela dit des deux erreurs~?
\end{exercise}
% ledger: ar6ch7:wv-lr

\begin{exercise}[$\star\star$]\label{exo:b2:anthropogenic-climate-change:7}
De 1850--1900 à 2006--2019, $\Delta N = \qty{0.59}{W.m^{-2}}$, $\Delta T = \qty{1.03}{\celsius}$
et $\Delta F = \qty{2.20}{W.m^{-2}}$. (a) Trouver $\lambda$ et la sensibilité effective. (b) Le
paramètre de rétroaction devrait devenir moins négatif de
$\lambda' = \qty{0.5}{W.m^{-2}.\celsius^{-1}}$ à mesure que le motif du réchauffement évolue.
Trouver la sensibilité $-F_{2\times}/(\lambda + \lambda')$. (c) Refaire (b) avec $\lambda' = 0.1$.
\end{exercise}
% ledger: ar6ch7:energy-budget, ar6ch7:pattern, ar6ch7:pattern-ecs

\begin{exercise}[$\star\star$]\label{exo:b2:anthropogenic-climate-change:8}
(a) À partir des \qty{0.617}{m} de hausse par degré de réchauffement de tout l'océan, de profondeur moyenne
\qty{3682}{m}, trouver le coefficient d'expansion effectif. (b) L'océan a gagné \qty{396}{ZJ} de 1971 à 2018.
Avec une efficacité d'expansion de \qty{0.113}{m} par \qty{e24}{J}, trouver l'élévation thermostérique, et
comparer avec les \qty{47.5}{mm} du bilan. (c) Une couche de \qty{700}{m} se réchauffe de \qty{0.2}{\celsius}
avec $\alpha = \qty{2e-4}{\celsius^{-1}}$~: de combien élève-t-elle le niveau de la mer~?
\end{exercise}
% ledger: ar6ch9:expansion, charette:ocean-volume, ar6ch7:inventory, ar6ch9:table95

\begin{exercise}[$\star\star$]\label{exo:b2:anthropogenic-climate-change:9}
Une empreinte modèle, sur quatre décennies, est $x = (0.10, 0.20, 0.30, 0.40)$ \unit{\celsius}~;
les observations sont $y = (0.12, 0.18, 0.33, 0.41)$ \unit{\celsius}. (a) Trouver le facteur
d'échelle $\beta$ des moindres carrés de $y = \beta x + \varepsilon$. (b) Son erreur standard,
estimée à partir des simulations sans forçage, vaut \num{0.08}. La réponse est-elle détectée~?
Est-elle compatible avec le modèle~? (c) Que signifierait un $\beta$ de 0.5 avec la même erreur~?
\end{exercise}

\begin{exercise}[$\star\star\star$]\label{exo:b2:anthropogenic-climate-change:10}
Deux empreintes, celle des gaz à effet de serre $g = (0.2, 0.4, 0.6, 0.8)$ et celle des aérosols $a = (-0.1, -0.2,
-0.2, -0.3)$, toutes deux en \unit{\celsius} sur quatre décennies, sont ajustées aux observations $y = (0.1,
0.2, 0.4, 0.5)$. (a) Écrire les équations normales de l'ajustement par moindres carrés $y = \beta_g g + \beta_a a
+ \varepsilon$ et les résoudre. (b) Trouver le réchauffement attribué à chaque facteur sur la dernière décennie. (c)
Pourquoi les incertitudes sur $\beta_g$ et $\beta_a$ augmentent-elles lorsque les deux empreintes se ressemblent~?
\end{exercise}

\begin{exercise}[$\star\star\star$]\label{exo:b2:anthropogenic-climate-change:11}
Appliquer la méthode du \omterm{def:b2:anthropogenic-climate-change:budget}{budget carbone restant} avec la TCRE de \qty{0.45}{\celsius} par \qty{1000}{Gt},
les \qty{210}{Gt} émises de 2015 à 2019 et \qty{26}{Gt} par degré pour les rétroactions du système
Terre. (a) Trouver le budget pour \qty{2}{\celsius}, soit \qty{0.93}{\celsius} au-dessus de 2010--2019,
si les autres facteurs ajoutent \qty{0.1}{\celsius}. (b) Quel réchauffement supplémentaire dû aux autres
facteurs rend ce budget égal aux \qty{1350}{Gt} évalués~? Est-il dans l'intervalle évalué~? (c) Montrer
que chaque \qty{0.1}{\celsius} de réchauffement hors CO$_2$ modifie tout budget de la même quantité.
\end{exercise}

\begin{exercise}[$\star\star\star$]\label{exo:b2:anthropogenic-climate-change:12}
La rétroaction nette a un intervalle très probable de $-1.81$ à \qty{-0.51}{W.m^{-2}.\celsius^{-1}}, autour
de $-1.16$. (a) Trouver l'ECS aux deux bornes et au centre. (b) Montrer qu'un intervalle de $\lambda$
symétrique autour de son centre donne un intervalle d'ECS qui ne l'est pas, et dire quel côté est le plus
long. (c) Pourquoi les évaluations trouvent-elles plus difficile d'exclure une ECS élevée qu'une ECS basse~?
\end{exercise}
% ledger: ar6ch7:table710

\begin{problem}[{Problème du week-end --- Lecture du budget}]\label{pb:b2:anthropogenic-climate-change:1}
Ce problème suit l'énergie et le carbone de l'évaluation, du forçage jusqu'au budget restant pour
\qty{1.5}{\celsius}. Utiliser les données des exercices, et~: ERF 1750--2019 (\unit{W.m^{-2}})~: dioxyde
de carbone 2.16, méthane 0.54, protoxyde d'azote 0.21, gaz halogénés 0.41, ozone 0.47, vapeur d'eau
stratosphérique 0.05, aérosols--rayonnement $-0.22$, aérosols--nuages $-0.84$, usage des sols $-0.20$,
carbone suie sur neige 0.08, traînées de condensation 0.06~; dioxyde de carbone \qty{410}{ppm} en 2019
et \qty{427.35}{ppm} en 2025~; réchauffement \qty{0.79}{W.m^{-2}} sur 2006--2018~; variations du bilan
énergétique $\Delta N = 0.59$, $\Delta F = \qty{2.20}{W.m^{-2}}$, $\Delta T = \qty{1.03}{\celsius}$~;
TCRE \qty{0.45}{\celsius} par \qty{1000}{Gt} de dioxyde de carbone (probable de 0.27 à 0.63)~;
\qty{2390}{Gt} émises de 1850 à 2019~; réchauffement humain \qty{1.07}{\celsius} pour 2010--2019~;
\qty{210}{Gt} émises de 2015 à 2019~; réchauffement hors CO$_2$ supplémentaire de \qty{0.1}{\celsius},
la borne basse de l'intervalle évalué de 0.1 à \qty{0.2}{\celsius}~; rétroactions du système Terre de
\qty{26}{Gt} par degré. Arrondir à trois chiffres significatifs.
% ledger: ipcc:erf-agents, ipcc:erf-2xco2, data:co2-mlo, ar6ch7:inventory, ar6ch7:energy-budget, ar6spm:tcre, ar6spm:budget, ar6ch5:budget-parts

\noindent\textbf{Partie I --- Le forçage.}
\begin{enumerate}
  \item Additionner le forçage des quatre gaz à effet de serre bien mélangés.
  \item Additionner le forçage des aérosols.
  \item Additionner tous les termes humains. Comparer avec les \qty{2.72}{W.m^{-2}} évalués.
  \item Vérifier le forçage du dioxyde de carbone en 2019 avec la loi logarithmique.
  \item Trouver le forçage du dioxyde de carbone en 2025 par la même loi.
  \item Quel terme a la plus grande incertitude, et pourquoi domine-t-il celle du total~?
\end{enumerate}

\noindent\textbf{Partie II --- Le bilan énergétique.}
\begin{enumerate}[resume]
  \item Quelle énergie le système climatique a-t-il gagnée de 2006 à 2018, en zettajoules~?
  \item Quelle part en est allée dans l'océan~?
  \item Trouver le paramètre de rétroaction à partir des variations du bilan énergétique.
  \item Trouver la sensibilité climatique effective à l'équilibre.
  \item Trouver la \omterm{def:b2:anthropogenic-climate-change:tcr}{réponse climatique transitoire} à partir des mêmes données.
  \item Trouver l'ECS à partir de la rétroaction nette de la figure, et comparer les deux valeurs.
  \item Donner la raison physique pour laquelle l'estimation historique est la plus basse.
\end{enumerate}

\noindent\textbf{Partie III --- Réchauffement par tonne.}
\begin{enumerate}[resume]
  \item Convertir la TCRE en degrés par \qty{1000}{PgC} de carbone.
  \item Convertir son intervalle probable.
  \item Convertir les émissions de 1850 à 2019 en pétagrammes de carbone.
  \item Quel réchauffement la TCRE donne-t-elle pour les émissions de 1850 à 2019~?
  \item Comparer avec le réchauffement humain de \qty{1.07}{\celsius}. Que dit la comparaison de l'effet
        net des autres facteurs humains~?
\end{enumerate}

\noindent\textbf{Partie IV --- Le budget restant pour \qty{1.5}{\celsius}.}
\begin{enumerate}[resume]
  \item Quel réchauffement reste-t-il au-dessus de 2010--2019~?
  \item Retrancher le réchauffement hors CO$_2$ et diviser par la TCRE~: quel budget reste à partir du milieu
        de 2010--2019~?
  \item Retrancher les émissions de 2015 à 2019.
  \item Retrancher les rétroactions du système Terre sur le réchauffement restant.
  \item De combien le budget change-t-il si les autres facteurs ajoutent plutôt \qty{0.2}{\celsius}~?
  \item En somme, quel est le \omterm{def:b2:anthropogenic-climate-change:budget}{budget carbone restant} pour \qty{1.5}{\celsius}, à partir du début de
        2020, calculé à partir de la TCRE~? Comparer avec les \qty{500}{Gt} évalués.
\end{enumerate}
\end{problem}
